% Kyriale — every chant of the repository in one A5 book: sprinkling rite,
% the Ordinary of Masses I, VIII, IX, XI and XVII, Credo I and III, Marian
% antiphons, the Requiem; contents, running heads and an index of chants.
% Build with `make` (lualatex, makeindex, lualatex twice; LilyPond in PATH).
\documentclass[11pt,twoside,openany]{book}

\usepackage[a5paper,top=11mm,bottom=13mm,inner=15mm,outer=11mm,
            includehead,headheight=14pt,headsep=3mm,footskip=7mm]{geometry}
\usepackage{fontspec}
\setmainfont{Libertinus Serif}
\usepackage[latin,main=polish]{babel}
\usepackage{microtype}
\usepackage{xcolor}
\usepackage{paracol}
\usepackage{titlesec}
\usepackage{titletoc}
\usepackage{fancyhdr}
\usepackage{makeidx}
\usepackage[
  program=lilypond,
  staffsize=14,
  includepaths={./,../shared/,../shared/scores/},
  pass-fonts,
  tmpdir=build-ly,
]{lyluatex}
\usepackage[hidelinks]{hyperref}

\newcommand{\tekstsize}{\footnotesize}
% Macros shared by the Mass booklets and the Kyriale; the including preamble
% loads fontspec, babel, xcolor and paracol and defines \tekstsize.
\definecolor{rubric}{RGB}{170,20,20}

% Parallel text: Latin left, Polish right, one row per sung phrase.
\newenvironment{tekst}
  {\par\vspace{0.6ex}\tekstsize
   \setcolumnwidth{0.48\textwidth,0.48\textwidth}%
   \setlength{\columnsep}{0.04\textwidth}%
   \begin{paracol}{2}}
  {\end{paracol}\par\vspace{0.5ex}}
% \w{Latin}{Polski}
\newcommand{\w}[2]{%
  \switchcolumn[0]*\foreignlanguage{latin}{#1}\par
  \switchcolumn #2\par\vspace{0.3ex}}
% Libertinus has no U+211F/U+2123: build the barred R and V.
\newcommand{\barred}[1]{#1\kern-.48em\raisebox{.04em}{/}\kern.06em}
\newcommand{\R}{\textcolor{rubric}{\barred{R}.}~}
\newcommand{\V}{\textcolor{rubric}{\barred{V}.}~}
\newcommand{\rubr}[1]{\textcolor{rubric}{\itshape #1}}

\setlength{\parindent}{0pt}
\setlength{\parskip}{0pt}
\makeindex

% --- running heads: part on the left page, Mass or chant on the right -----
\pagestyle{fancy}
\fancyhf{}
% Heads come from macros, not marks (paracol upsets marks); both change only
% right after a \clearpage, so a page never shows the next section's name.
\newcommand{\headleft}{}
\newcommand{\headright}{}
\fancyhead[LE]{\footnotesize\scshape\headleft}
\fancyhead[RO]{\footnotesize\scshape\headright}
\fancyfoot[C]{\footnotesize\thepage}
\renewcommand{\headrulewidth}{0.3pt}
% (plain is left as the class defines it: redefining it here wipes the fancy
% heads once a page uses it)
\addto\captionspolish{\renewcommand{\indexname}{Indeks}}

% --- headings --------------------------------------------------------------
% \dzial{Latin title}{Polish subtitle}: one of the four parts of the book.
\newcommand{\currentmass}{}
\newcommand{\masskey}{}
\newif\ifdzialtop  % a part heading was just set: the first Mass shares its page
\newcommand{\dzial}[2]{%
  \clearpage
  \thispagestyle{plain}%
  \dzialtoptrue
  \renewcommand{\currentmass}{}%
  \gdef\headleft{#1}\gdef\headright{#1}%
  \phantomsection\addcontentsline{toc}{part}{#1}%
  \begin{center}
    {\LARGE\scshape #1}\par\smallskip
    {\small\itshape #2}\par
  \end{center}
  \vspace{1ex}}
% \missa{number}{name}{sort key}{occasion}: a Mass within the Ordinarium.
\newcommand{\missa}[4]{%
  \ifdzialtop\dzialtopfalse\else\clearpage\fi
  \renewcommand{\currentmass}{Missa #1}%
  \renewcommand{\masskey}{#3}%
  \phantomsection
  \if\relax\detokenize{#2}\relax
    \gdef\headright{Missa #1}%
    \addcontentsline{toc}{chapter}{Missa #1}%
    \index{Missa #3@Missa #1}%
  \else
    \gdef\headright{Missa #1 · #2}%
    \addcontentsline{toc}{chapter}{Missa #1 \textit{#2}}%
    \index{Missa #3@Missa #1 \textit{#2}}%
    \index{#2|see{Missa #1}}%
  \fi
  \begin{center}
    {\Large\scshape Missa #1}\par
    \if\relax\detokenize{#2}\relax\else{\large\itshape #2}\par\fi
    \if\relax\detokenize{#4}\relax\else{\small\itshape #4}\par\fi
  \end{center}
  \vspace{0.5ex}}
% Contents and index titles, in the style of the part headings.
\titleformat{\chapter}[block]{\centering\LARGE\scshape}{}{0pt}{}
\titlespacing*{\chapter}{0pt}{0pt}{3ex}
% \chant{Latin title}{mode}: one chant; listed in the contents and the index.
\titleformat{\section}{\centering\large\scshape}{}{0pt}{}
\titlespacing*{\section}{0pt}{1.6ex plus 1ex}{0.2ex}
\newcommand{\chant}[2]{%
  \section*{\foreignlanguage{latin}{#1}}%
  \phantomsection\addcontentsline{toc}{section}{#1}%
  \ifx\currentmass\empty
    \index{#1}%
  \else
    \index{#1!Missa \masskey@\currentmass}%
  \fi
  {\centering\small\itshape #2\par}%
  \vspace{0.6ex}}

% Credo is printed once, in its own part; the Masses point to it.
\newcommand{\credo}[1]{%
  \par\nopagebreak[4]\vspace{1ex}\nopagebreak[4]
  {\centering\small\rubr{Credo \MakeUppercase{#1} — zob. s.~\pageref{credo-#1}.}\par}}
\newcommand{\antiphonsheading}{}

% --- contents --------------------------------------------------------------
\setcounter{tocdepth}{1}
\titlecontents{part}[0pt]{\addvspace{1.2ex}\scshape}{}{}{\hfill\contentspage}
\titlecontents{chapter}[1em]{\addvspace{0.4ex}}{}{}{\titlerule*[0.6em]{.}\contentspage}
\titlecontents{section}[2.2em]{\small}{}{}{\titlerule*[0.6em]{.}\contentspage}

\begin{document}
\frontmatter
% --- title page --------------------------------------------------------------
\begin{titlepage}
  \centering
  \vspace*{20mm}
  {\fontsize{38}{40}\selectfont\scshape Kyriale}\par\vspace{4mm}
  {\Large\itshape Śpiewy stałe Mszy świętej}\par\vspace{14mm}
  {\small Asperges me · Vidi aquam\\[0.4ex]
   Missa I · VIII · IX · XI · XVII\\[0.4ex]
   Credo I · III\\[0.4ex]
   Antyfony maryjne\\[0.4ex]
   Missa pro defunctis\par}
  \vfill
  {\footnotesize Kyriale Romanum · zapis współczesny wg Vatican II Hymnal (2011)\\
   tłumaczenie polskie wg Mszału rzymskiego dla diecezji polskich\par}
\end{titlepage}

% --- verso of the title page -------------------------------------------------
\thispagestyle{empty}
\vspace*{\fill}
{\footnotesize
Zapis nutowy na podstawie GregoBase (gregobase.selapa.net); wysokość, łuki
i podziały wg Vatican II Hymnal (Corpus Christi Watershed, 2011). Źródła
i transkrybenci: plik ATTRIBUTION.md w repozytorium.\par\medskip
Opracowanie przekazane do domeny publicznej (CC0 1.0). Teksty polskie wg
\textit{Mszału rzymskiego dla diecezji polskich}.\par}
\clearpage

\gdef\headleft{Spis treści}\gdef\headright{Spis treści}
\tableofcontents
\clearpage

\mainmatter
% --- I. sprinkling rite ------------------------------------------------------
\dzial{Ritus aspersionis}{Obrzęd pokropienia wodą święconą}
% Sprinkling rite, sung before Mass; \input by every booklet.
% ---------------------------------------------------------------------------
\chant{Asperges me}{tonus VII}
\lilypondfile{../shared/scores/asperges.ly}
\begin{tekst}
\w{Aspérges me, Dómine, hyssópo, et mundábor:}
  {Pokrop mnie Panie hyzopem a stanę się czysty,}
\w{lavábis me, et super nivem dealbábor.}
  {obmyj mnie, a nad śnieg wybieleję.}
\w{\rubr{Ps. 50.} Miserére mei, Deus, secúndum magnam misericórdiam tuam.}
  {\rubr{Ps. 51.} Zmiłuj się nade mną, Boże, według wielkiego miłosierdzia Twego.}
\w{Glória Patri, et Fílio, et Spirítui Sancto:}
  {Chwała Ojcu, i Synowi, i Duchowi Świętemu,}
\w{Sicut erat in princípio, et nunc, et semper, et in saécula saeculórum. Amen.}
  {jak była na początku, teraz i zawsze i na wieki wieków. Amen.}
\w{\rubr{Powtarza się antyfonę} Aspérges me.}
  {\rubr{Powtarza się antyfonę} Pokrop mnie.}
\end{tekst}


% ---------------------------------------------------------------------------
\chant{Asperges me II}{tonus IV}
{\centering\small\rubr{Melodia do wyboru (ad libitum) zamiast powyższej.}\par}
\lilypondfile{../shared/scores/asperges-ii.ly}
{\centering\small\rubr{Psalm, Glória Patri i powtórzenie antyfony jak wyżej.}\par}

\clearpage
% ---------------------------------------------------------------------------
\chant{Vidi aquam}{tonus VIII}
{\centering\small\rubr{W okresie wielkanocnym, od Niedzieli Zmartwychwstania do Zesłania
Ducha Świętego, zamiast Asperges me śpiewa się:}\par}
\lilypondfile{../shared/scores/vidi-aquam.ly}
\begin{tekst}
\w{Vidi aquam egrediéntem de templo, a látere dextro, allelúja:}
  {Widziałem wodę wypływającą z prawej strony świątyni, alleluja,}
\w{et omnes, ad quos pervénit aqua ista, salvi facti sunt, et dicent,
   allelúja, allelúja.}
  {a wszyscy, do których ta woda dotarła, zostali zbawieni i śpiewają:
   alleluja, alleluja.}
\w{\rubr{Ps. 117.} Confitémini Dómino quóniam bonus:
   quóniam in saéculum misericórdia ejus.}
  {\rubr{Ps. 118.} Wysławiajcie Pana, bo jest dobry,
   bo Jego miłosierdzie trwa na wieki.}
\w{Glória Patri, et Fílio, et Spirítui Sancto:}
  {Chwała Ojcu, i Synowi, i Duchowi Świętemu,}
\w{Sicut erat in princípio, et nunc, et semper, et in saécula saeculórum. Amen.}
  {jak była na początku, teraz i zawsze i na wieki wieków. Amen.}
\w{\rubr{Powtarza się antyfonę} Vidi aquam.}
  {\rubr{Powtarza się antyfonę} Widziałem wodę.}
\end{tekst}


% --- II. Ordinary of the Mass -------------------------------------------------
\dzial{Ordinarium Missae}{Części stałe Mszy świętej}
\missa{I}{Lux et origo}{01}{Tempore paschali · w okresie wielkanocnym}
% Ordinary of Missa I: \input by missa-i.tex and ../kyriale/kyriale.tex.
% Score paths are relative to a sibling directory, so they work from both.
% ---------------------------------------------------------------------------
\chant{Kyrie}{tonus VIII}
\lilypondfile{../missa-i/scores/kyrie.ly}
\begin{tekst}
\w{Kýrie, eléison. Christe, eléison. Kýrie, eléison.}
  {Panie, zmiłuj się nad nami. Chryste, zmiłuj się nad nami.
   Panie, zmiłuj się nad nami.}
\end{tekst}

\clearpage
% ---------------------------------------------------------------------------
\chant{Gloria}{tonus IV}
{\centering\small\rubr{Jeśli to możliwe, intonuje sam celebrans.}\par}
\lilypondfile{../missa-i/scores/gloria.ly}
\begin{tekst}
\w{Glória in excélsis Deo.}{Chwała na wysokości Bogu,}
\w{Et in terra pax homínibus bonae voluntátis.}{a na ziemi pokój ludziom dobrej woli.}
\w{Laudámus te. Benedícimus te. Adorámus te. Glorificámus te.}
  {Chwalimy Cię. Błogosławimy Cię. Wielbimy Cię. Wysławiamy Cię.}
\w{Grátias ágimus tibi propter magnam glóriam tuam.}
  {Dzięki Ci składamy, bo wielka jest chwała Twoja.}
\w{Dómine Deus, Rex caeléstis, Deus Pater omnípotens.}
  {Panie Boże, Królu nieba, Boże, Ojcze wszechmogący.}
\w{Dómine Fili unigénite, Jesu Christe.}
  {Panie, Synu Jednorodzony, Jezu Chryste.}
\w{Dómine Deus, Agnus Dei, Fílius Patris.}
  {Panie Boże, Baranku Boży, Synu Ojca.}
\w{Qui tollis peccáta mundi, miserére nobis.}
  {Który gładzisz grzechy świata, zmiłuj się nad nami.}
\w{Qui tollis peccáta mundi, súscipe deprecatiónem nostram.}
  {Który gładzisz grzechy świata, przyjm błaganie nasze.}
\w{Qui sedes ad déxteram Patris, miserére nobis.}
  {Który siedzisz po prawicy Ojca, zmiłuj się nad nami.}
\w{Quóniam tu solus sanctus. Tu solus Dóminus.}
  {Albowiem tylko Tyś jest święty. Tylko Tyś jest Panem.}
\w{Tu solus Altíssimus, Jesu Christe.}
  {Tylko Tyś Najwyższy, Jezu Chryste.}
\w{Cum Sancto Spíritu, in glória Dei Patris. Amen.}
  {Z Duchem Świętym w chwale Boga Ojca. Amen.}
\end{tekst}

\credo{iii}

\clearpage
% ---------------------------------------------------------------------------
\chant{Sanctus}{tonus IV}
\lilypondfile{../missa-i/scores/sanctus.ly}
\begin{tekst}
\w{Sanctus, Sanctus, Sanctus Dóminus Deus Sábaoth.}
  {Święty, Święty, Święty, Pan Bóg Zastępów.}
\w{Pleni sunt caeli et terra glória tua.}
  {Pełne są niebiosa i ziemia chwały Twojej.}
\w{Hosánna in excélsis.}{Hosanna na wysokości.}
\w{Benedíctus qui venit in nómine Dómini.}
  {Błogosławiony, który idzie w imię Pańskie.}
\w{Hosánna in excélsis.}{Hosanna na wysokości.}
\end{tekst}

% ---------------------------------------------------------------------------
\chant{Agnus Dei}{tonus IV}
\lilypondfile{../missa-i/scores/agnus.ly}
\begin{tekst}
\w{Agnus Dei, qui tollis peccáta mundi: miserére nobis.}
  {Baranku Boży, który gładzisz grzechy świata, zmiłuj się nad nami.}
\w{Agnus Dei, qui tollis peccáta mundi: miserére nobis.}
  {Baranku Boży, który gładzisz grzechy świata, zmiłuj się nad nami.}
\w{Agnus Dei, qui tollis peccáta mundi: dona nobis pacem.}
  {Baranku Boży, który gładzisz grzechy świata, obdarz nas pokojem.}
\end{tekst}

\clearpage
% ---------------------------------------------------------------------------
\chant{Ite, missa est}{tonus VII}
\lilypondfile{../missa-i/scores/ite.ly}
\begin{tekst}
\w{\V Ite, missa est.}{\V Idźcie w pokoju Chrystusa.}
\w{\R Deo grátias.}{\R Bogu niech będą dzięki.}
\end{tekst}

\par\medskip
{\small\itshape\textcolor{rubric}{W oktawie Wielkanocy i w uroczystość
Zesłania Ducha Świętego:}\par}
\lilypondfile{../missa-i/scores/ite-paschale.ly}
\begin{tekst}
\w{\V Ite, missa est, allelúia, allelúia.}
  {\V Idźcie w pokoju Chrystusa, alleluja, alleluja.}
\w{\R Deo grátias, allelúia, allelúia.}
  {\R Bogu niech będą dzięki, alleluja, alleluja.}
\end{tekst}

\missa{VIII}{De Angelis}{08}{}
% Ordinary of Missa VIII: \input by missa-viii.tex and ../kyriale/kyriale.tex.
% Score paths are relative to a sibling directory, so they work from both.
% ---------------------------------------------------------------------------
\chant{Kyrie}{tonus V}
\lilypondfile{../missa-viii/scores/kyrie.ly}
\par\smallskip
{\small\itshape\textcolor{rubric}{Jeśli Kyrie śpiewa się w tradycyjnej formie
dziewięciokrotnej, jako dziewiąte (ostatnie) wezwanie śpiewa się:}\par}
\lilypondfile{../missa-viii/scores/kyrie-ix.ly}
\begin{tekst}
\w{Kýrie, eléison. Christe, eléison. Kýrie, eléison.}
  {Panie, zmiłuj się nad nami. Chryste, zmiłuj się nad nami.
   Panie, zmiłuj się nad nami.}
\end{tekst}

\clearpage
% ---------------------------------------------------------------------------
\chant{Gloria}{tonus V}
{\centering\small\rubr{Jeśli to możliwe, intonuje sam celebrans.}\par}
\lilypondfile{../missa-viii/scores/gloria.ly}
\begin{tekst}
\w{Glória in excélsis Deo.}{Chwała na wysokości Bogu,}
\w{Et in terra pax homínibus bonae voluntátis.}{a na ziemi pokój ludziom dobrej woli.}
\w{Laudámus te. Benedícimus te. Adorámus te. Glorificámus te.}
  {Chwalimy Cię. Błogosławimy Cię. Wielbimy Cię. Wysławiamy Cię.}
\w{Grátias ágimus tibi propter magnam glóriam tuam.}
  {Dzięki Ci składamy, bo wielka jest chwała Twoja.}
\w{Dómine Deus, Rex caeléstis, Deus Pater omnípotens.}
  {Panie Boże, Królu nieba, Boże, Ojcze wszechmogący.}
\w{Dómine Fili unigénite, Jesu Christe.}
  {Panie, Synu Jednorodzony, Jezu Chryste.}
\w{Dómine Deus, Agnus Dei, Fílius Patris.}
  {Panie Boże, Baranku Boży, Synu Ojca.}
\w{Qui tollis peccáta mundi, miserére nobis.}
  {Który gładzisz grzechy świata, zmiłuj się nad nami.}
\w{Qui tollis peccáta mundi, súscipe deprecatiónem nostram.}
  {Który gładzisz grzechy świata, przyjm błaganie nasze.}
\w{Qui sedes ad déxteram Patris, miserére nobis.}
  {Który siedzisz po prawicy Ojca, zmiłuj się nad nami.}
\w{Quóniam tu solus sanctus. Tu solus Dóminus.}
  {Albowiem tylko Tyś jest święty. Tylko Tyś jest Panem.}
\w{Tu solus Altíssimus, Jesu Christe.}
  {Tylko Tyś Najwyższy, Jezu Chryste.}
\w{Cum Sancto Spíritu, in glória Dei Patris. Amen.}
  {Z Duchem Świętym w chwale Boga Ojca. Amen.}
\end{tekst}

\credo{iii}

\clearpage
% ---------------------------------------------------------------------------
\chant{Sanctus}{tonus VI}
\lilypondfile{../missa-viii/scores/sanctus.ly}
\begin{tekst}
\w{Sanctus, Sanctus, Sanctus Dóminus Deus Sábaoth.}
  {Święty, Święty, Święty, Pan Bóg Zastępów.}
\w{Pleni sunt caeli et terra glória tua.}
  {Pełne są niebiosa i ziemia chwały Twojej.}
\w{Hosánna in excélsis.}{Hosanna na wysokości.}
\w{Benedíctus qui venit in nómine Dómini.}
  {Błogosławiony, który idzie w imię Pańskie.}
\w{Hosánna in excélsis.}{Hosanna na wysokości.}
\end{tekst}

% ---------------------------------------------------------------------------
\chant{Agnus Dei}{tonus VI}
\lilypondfile{../missa-viii/scores/agnus.ly}
\begin{tekst}
\w{Agnus Dei, qui tollis peccáta mundi: miserére nobis.}
  {Baranku Boży, który gładzisz grzechy świata, zmiłuj się nad nami.}
\w{Agnus Dei, qui tollis peccáta mundi: miserére nobis.}
  {Baranku Boży, który gładzisz grzechy świata, zmiłuj się nad nami.}
\w{Agnus Dei, qui tollis peccáta mundi: dona nobis pacem.}
  {Baranku Boży, który gładzisz grzechy świata, obdarz nas pokojem.}
\end{tekst}

% ---------------------------------------------------------------------------
\chant{Ite, missa est}{tonus V}
\lilypondfile{../missa-viii/scores/ite.ly}
\begin{tekst}
\w{\V Ite, missa est.}{\V Idźcie w pokoju Chrystusa.}
\w{\R Deo grátias.}{\R Bogu niech będą dzięki.}
\end{tekst}

\missa{IX}{Cum jubilo}{09}{In festis B. Mariae Virginis · w święta Najświętszej Maryi Panny}
% Ordinary of Missa IX: \input by missa-ix.tex and ../kyriale/kyriale.tex.
% Score paths are relative to a sibling directory, so they work from both.
% ---------------------------------------------------------------------------
\chant{Kyrie}{tonus I}
\lilypondfile{../missa-ix/scores/kyrie.ly}
\begin{tekst}
\w{Kýrie, eléison. Christe, eléison. Kýrie, eléison.}
  {Panie, zmiłuj się nad nami. Chryste, zmiłuj się nad nami.
   Panie, zmiłuj się nad nami.}
\end{tekst}

\clearpage
% ---------------------------------------------------------------------------
\chant{Gloria}{tonus VII}
{\centering\small\rubr{Jeśli to możliwe, intonuje sam celebrans.}\par}
\lilypondfile{../missa-ix/scores/gloria.ly}
\begin{tekst}
\w{Glória in excélsis Deo.}{Chwała na wysokości Bogu,}
\w{Et in terra pax homínibus bonae voluntátis.}{a na ziemi pokój ludziom dobrej woli.}
\w{Laudámus te. Benedícimus te. Adorámus te. Glorificámus te.}
  {Chwalimy Cię. Błogosławimy Cię. Wielbimy Cię. Wysławiamy Cię.}
\w{Grátias ágimus tibi propter magnam glóriam tuam.}
  {Dzięki Ci składamy, bo wielka jest chwała Twoja.}
\w{Dómine Deus, Rex caeléstis, Deus Pater omnípotens.}
  {Panie Boże, Królu nieba, Boże, Ojcze wszechmogący.}
\w{Dómine Fili unigénite, Jesu Christe.}
  {Panie, Synu Jednorodzony, Jezu Chryste.}
\w{Dómine Deus, Agnus Dei, Fílius Patris.}
  {Panie Boże, Baranku Boży, Synu Ojca.}
\w{Qui tollis peccáta mundi, miserére nobis.}
  {Który gładzisz grzechy świata, zmiłuj się nad nami.}
\w{Qui tollis peccáta mundi, súscipe deprecatiónem nostram.}
  {Który gładzisz grzechy świata, przyjm błaganie nasze.}
\w{Qui sedes ad déxteram Patris, miserére nobis.}
  {Który siedzisz po prawicy Ojca, zmiłuj się nad nami.}
\w{Quóniam tu solus sanctus. Tu solus Dóminus.}
  {Albowiem tylko Tyś jest święty. Tylko Tyś jest Panem.}
\w{Tu solus Altíssimus, Jesu Christe.}
  {Tylko Tyś Najwyższy, Jezu Chryste.}
\w{Cum Sancto Spíritu, in glória Dei Patris. Amen.}
  {Z Duchem Świętym w chwale Boga Ojca. Amen.}
\end{tekst}

\credo{iii}

\clearpage
% ---------------------------------------------------------------------------
\chant{Sanctus}{tonus V}
\lilypondfile{../missa-ix/scores/sanctus.ly}
\begin{tekst}
\w{Sanctus, Sanctus, Sanctus Dóminus Deus Sábaoth.}
  {Święty, Święty, Święty, Pan Bóg Zastępów.}
\w{Pleni sunt caeli et terra glória tua.}
  {Pełne są niebiosa i ziemia chwały Twojej.}
\w{Hosánna in excélsis.}{Hosanna na wysokości.}
\w{Benedíctus qui venit in nómine Dómini.}
  {Błogosławiony, który idzie w imię Pańskie.}
\w{Hosánna in excélsis.}{Hosanna na wysokości.}
\end{tekst}

% ---------------------------------------------------------------------------
\chant{Agnus Dei}{tonus V}
\lilypondfile{../missa-ix/scores/agnus.ly}
\begin{tekst}
\w{Agnus Dei, qui tollis peccáta mundi: miserére nobis.}
  {Baranku Boży, który gładzisz grzechy świata, zmiłuj się nad nami.}
\w{Agnus Dei, qui tollis peccáta mundi: miserére nobis.}
  {Baranku Boży, który gładzisz grzechy świata, zmiłuj się nad nami.}
\w{Agnus Dei, qui tollis peccáta mundi: dona nobis pacem.}
  {Baranku Boży, który gładzisz grzechy świata, obdarz nas pokojem.}
\end{tekst}

% ---------------------------------------------------------------------------
\chant{Ite, missa est}{tonus I}
\lilypondfile{../missa-ix/scores/ite.ly}
\begin{tekst}
\w{\V Ite, missa est.}{\V Idźcie w pokoju Chrystusa.}
\w{\R Deo grátias.}{\R Bogu niech będą dzięki.}
\end{tekst}

\missa{XI}{Orbis factor}{11}{In dominicis per annum · w niedziele w ciągu roku}
% Ordinary of Missa XI: \input by missa-xi.tex and ../kyriale/kyriale.tex.
% Score paths are relative to a sibling directory, so they work from both.
% ---------------------------------------------------------------------------
\chant{Kyrie}{tonus I}
\lilypondfile{../missa-xi/scores/kyrie.ly}
\begin{tekst}
\w{Kýrie, eléison. Christe, eléison. Kýrie, eléison.}
  {Panie, zmiłuj się nad nami. Chryste, zmiłuj się nad nami.
   Panie, zmiłuj się nad nami.}
\end{tekst}

\clearpage
% ---------------------------------------------------------------------------
\chant{Gloria}{tonus II}
{\centering\small\rubr{Jeśli to możliwe, intonuje sam celebrans.}\par}
\lilypondfile{../missa-xi/scores/gloria.ly}
\begin{tekst}
\w{Glória in excélsis Deo.}{Chwała na wysokości Bogu,}
\w{Et in terra pax homínibus bonae voluntátis.}{a na ziemi pokój ludziom dobrej woli.}
\w{Laudámus te. Benedícimus te. Adorámus te. Glorificámus te.}
  {Chwalimy Cię. Błogosławimy Cię. Wielbimy Cię. Wysławiamy Cię.}
\w{Grátias ágimus tibi propter magnam glóriam tuam.}
  {Dzięki Ci składamy, bo wielka jest chwała Twoja.}
\w{Dómine Deus, Rex caeléstis, Deus Pater omnípotens.}
  {Panie Boże, Królu nieba, Boże, Ojcze wszechmogący.}
\w{Dómine Fili unigénite, Jesu Christe.}
  {Panie, Synu Jednorodzony, Jezu Chryste.}
\w{Dómine Deus, Agnus Dei, Fílius Patris.}
  {Panie Boże, Baranku Boży, Synu Ojca.}
\w{Qui tollis peccáta mundi, miserére nobis.}
  {Który gładzisz grzechy świata, zmiłuj się nad nami.}
\w{Qui tollis peccáta mundi, súscipe deprecatiónem nostram.}
  {Który gładzisz grzechy świata, przyjm błaganie nasze.}
\w{Qui sedes ad déxteram Patris, miserére nobis.}
  {Który siedzisz po prawicy Ojca, zmiłuj się nad nami.}
\w{Quóniam tu solus sanctus. Tu solus Dóminus.}
  {Albowiem tylko Tyś jest święty. Tylko Tyś jest Panem.}
\w{Tu solus Altíssimus, Jesu Christe.}
  {Tylko Tyś Najwyższy, Jezu Chryste.}
\w{Cum Sancto Spíritu, in glória Dei Patris. Amen.}
  {Z Duchem Świętym w chwale Boga Ojca. Amen.}
\end{tekst}

\credo{i}

\clearpage
% ---------------------------------------------------------------------------
\chant{Sanctus}{tonus II}
\lilypondfile{../missa-xi/scores/sanctus.ly}
\begin{tekst}
\w{Sanctus, Sanctus, Sanctus Dóminus Deus Sábaoth.}
  {Święty, Święty, Święty, Pan Bóg Zastępów.}
\w{Pleni sunt caeli et terra glória tua.}
  {Pełne są niebiosa i ziemia chwały Twojej.}
\w{Hosánna in excélsis.}{Hosanna na wysokości.}
\w{Benedíctus qui venit in nómine Dómini.}
  {Błogosławiony, który idzie w imię Pańskie.}
\w{Hosánna in excélsis.}{Hosanna na wysokości.}
\end{tekst}

% ---------------------------------------------------------------------------
\chant{Agnus Dei}{tonus I}
\lilypondfile{../missa-xi/scores/agnus.ly}
\begin{tekst}
\w{Agnus Dei, qui tollis peccáta mundi: miserére nobis.}
  {Baranku Boży, który gładzisz grzechy świata, zmiłuj się nad nami.}
\w{Agnus Dei, qui tollis peccáta mundi: miserére nobis.}
  {Baranku Boży, który gładzisz grzechy świata, zmiłuj się nad nami.}
\w{Agnus Dei, qui tollis peccáta mundi: dona nobis pacem.}
  {Baranku Boży, który gładzisz grzechy świata, obdarz nas pokojem.}
\end{tekst}

% ---------------------------------------------------------------------------
\chant{Ite, missa est}{tonus I}
\lilypondfile{../missa-xi/scores/ite.ly}
\begin{tekst}
\w{\V Ite, missa est.}{\V Idźcie w pokoju Chrystusa.}
\w{\R Deo grátias.}{\R Bogu niech będą dzięki.}
\end{tekst}

\missa{XVII}{}{17}{In dominicis Adventus et Quadragesimae · w niedziele Adwentu i Wielkiego Postu}
% Ordinary of Missa XVII: \input by missa-xvii.tex and ../kyriale/kyriale.tex.
% Score paths are relative to a sibling directory, so they work from both.
% ---------------------------------------------------------------------------
\chant{Kyrie}{tonus I}
\lilypondfile{../missa-xvii/scores/kyrie.ly}
\begin{tekst}
\w{Kýrie, eléison. Christe, eléison. Kýrie, eléison.}
  {Panie, zmiłuj się nad nami. Chryste, zmiłuj się nad nami.
   Panie, zmiłuj się nad nami.}
\end{tekst}

{\centering\small\rubr{W Adwencie i Wielkim Poście nie śpiewa się Gloria.}\par}

\credo{i}

\clearpage
% ---------------------------------------------------------------------------
\chant{Sanctus}{tonus V}
\lilypondfile{../missa-xvii/scores/sanctus.ly}
\begin{tekst}
\w{Sanctus, Sanctus, Sanctus Dóminus Deus Sábaoth.}
  {Święty, Święty, Święty, Pan Bóg Zastępów.}
\w{Pleni sunt caeli et terra glória tua.}
  {Pełne są niebiosa i ziemia chwały Twojej.}
\w{Hosánna in excélsis.}{Hosanna na wysokości.}
\w{Benedíctus qui venit in nómine Dómini.}
  {Błogosławiony, który idzie w imię Pańskie.}
\w{Hosánna in excélsis.}{Hosanna na wysokości.}
\end{tekst}

% ---------------------------------------------------------------------------
\chant{Agnus Dei}{tonus V}
\lilypondfile{../missa-xvii/scores/agnus.ly}
\begin{tekst}
\w{Agnus Dei, qui tollis peccáta mundi: miserére nobis.}
  {Baranku Boży, który gładzisz grzechy świata, zmiłuj się nad nami.}
\w{Agnus Dei, qui tollis peccáta mundi: miserére nobis.}
  {Baranku Boży, który gładzisz grzechy świata, zmiłuj się nad nami.}
\w{Agnus Dei, qui tollis peccáta mundi: dona nobis pacem.}
  {Baranku Boży, który gładzisz grzechy świata, obdarz nas pokojem.}
\end{tekst}

% ---------------------------------------------------------------------------
\chant{Ite, missa est}{tonus IV}
\lilypondfile{../missa-xvii/scores/ite.ly}
\begin{tekst}
\w{\V Ite, missa est.}{\V Idźcie w pokoju Chrystusa.}
\w{\R Deo grátias.}{\R Bogu niech będą dzięki.}
\end{tekst}


% --- III. Credo ---------------------------------------------------------------
\dzial{Credo}{Wyznanie wiary}
\label{credo-i}
% Credo I; \input by the booklets that use it.
% ---------------------------------------------------------------------------
\chant{Credo I}{tonus IV}
{\centering\small\rubr{Jeśli to możliwe, intonuje sam celebrans.}\par}
\lilypondfile{../shared/scores/credo-i.ly}
% the score fills its page; keep the text together on the next one
\clearpage
\begin{tekst}
\w{Credo in unum Deum,}{Wierzę w jednego Boga,}
\w{Patrem omnipoténtem, factórem caeli et terrae,
   visibílium ómnium et invisibílium.}
  {Ojca wszechmogącego, Stworzyciela nieba i ziemi,
   wszystkich rzeczy widzialnych i niewidzialnych.}
\w{Et in unum Dóminum Jesum Christum, Fílium Dei unigénitum.}
  {I w jednego Pana Jezusa Chrystusa, Syna Bożego Jednorodzonego,}
\w{Et ex Patre natum ante ómnia saécula.}
  {który z Ojca jest zrodzony przed wszystkimi wiekami.}
\w{Deum de Deo, lumen de lúmine, Deum verum de Deo vero.}
  {Bóg z Boga, Światłość ze Światłości, Bóg prawdziwy z Boga prawdziwego.}
\w{Génitum, non factum, consubstantiálem Patri: per quem ómnia facta sunt.}
  {Zrodzony, a nie stworzony, współistotny Ojcu, a przez Niego wszystko się stało.}
\w{Qui propter nos hómines, et propter nostram salútem descéndit de caelis.}
  {On to dla nas, ludzi, i dla naszego zbawienia zstąpił z nieba.}
\w{Et incarnátus est de Spíritu Sancto ex María Vírgine: Et homo factus est.}
  {I za sprawą Ducha Świętego przyjął ciało z Maryi Dziewicy, i stał się człowiekiem.}
\w{Crucifíxus étiam pro nobis: sub Póntio Piláto passus, et sepúltus est.}
  {Ukrzyżowany również za nas, pod Poncjuszem Piłatem został umęczony i pogrzebany.}
\w{Et resurréxit tértia die, secúndum Scriptúras.}
  {I zmartwychwstał dnia trzeciego, jak oznajmia Pismo.}
\w{Et ascéndit in caelum: sedet ad déxteram Patris.}
  {I wstąpił do nieba; siedzi po prawicy Ojca.}
\w{Et íterum ventúrus est cum glória, judicáre vivos et mórtuos:
   cujus regni non erit finis.}
  {I powtórnie przyjdzie w chwale sądzić żywych i umarłych,
   a królestwu Jego nie będzie końca.}
\w{Et in Spíritum Sanctum, Dóminum, et vivificántem:
   qui ex Patre Filióque procédit.}
  {Wierzę w Ducha Świętego, Pana i Ożywiciela,
   który od Ojca i Syna pochodzi.}
\w{Qui cum Patre et Fílio simul adorátur, et conglorificátur:
   qui locútus est per Prophétas.}
  {Który z Ojcem i Synem wspólnie odbiera uwielbienie i chwałę;
   który mówił przez Proroków.}
\w{Et unam sanctam cathólicam et apostólicam Ecclésiam.}
  {Wierzę w jeden, święty, powszechny i apostolski Kościół.}
\w{Confíteor unum baptísma in remissiónem peccatórum.}
  {Wyznaję jeden chrzest na odpuszczenie grzechów.}
\w{Et exspécto resurrectiónem mortuórum.}
  {I oczekuję wskrzeszenia umarłych.}
\w{Et vitam ventúri saéculi. Amen.}
  {I życia wiecznego w przyszłym świecie. Amen.}
\end{tekst}


\clearpage
\label{credo-iii}
% Credo III; \input by the booklets that use it.
% ---------------------------------------------------------------------------
\chant{Credo III}{tonus V}
{\centering\small\rubr{Jeśli to możliwe, intonuje sam celebrans.}\par}
\lilypondfile{../shared/scores/credo-iii.ly}
% the score fills its page; keep the text together on the next one
\clearpage
\begin{tekst}
\w{Credo in unum Deum,}{Wierzę w jednego Boga,}
\w{Patrem omnipoténtem, factórem caeli et terrae,
   visibílium ómnium et invisibílium.}
  {Ojca wszechmogącego, Stworzyciela nieba i ziemi,
   wszystkich rzeczy widzialnych i niewidzialnych.}
\w{Et in unum Dóminum Jesum Christum, Fílium Dei unigénitum.}
  {I w jednego Pana Jezusa Chrystusa, Syna Bożego Jednorodzonego,}
\w{Et ex Patre natum ante ómnia saécula.}
  {który z Ojca jest zrodzony przed wszystkimi wiekami.}
\w{Deum de Deo, lumen de lúmine, Deum verum de Deo vero.}
  {Bóg z Boga, Światłość ze Światłości, Bóg prawdziwy z Boga prawdziwego.}
\w{Génitum, non factum, consubstantiálem Patri: per quem ómnia facta sunt.}
  {Zrodzony, a nie stworzony, współistotny Ojcu, a przez Niego wszystko się stało.}
\w{Qui propter nos hómines, et propter nostram salútem descéndit de caelis.}
  {On to dla nas, ludzi, i dla naszego zbawienia zstąpił z nieba.}
\w{Et incarnátus est de Spíritu Sancto ex María Vírgine: Et homo factus est.}
  {I za sprawą Ducha Świętego przyjął ciało z Maryi Dziewicy, i stał się człowiekiem.}
\w{Crucifíxus étiam pro nobis: sub Póntio Piláto passus, et sepúltus est.}
  {Ukrzyżowany również za nas, pod Poncjuszem Piłatem został umęczony i pogrzebany.}
\w{Et resurréxit tértia die, secúndum Scriptúras.}
  {I zmartwychwstał dnia trzeciego, jak oznajmia Pismo.}
\w{Et ascéndit in caelum: sedet ad déxteram Patris.}
  {I wstąpił do nieba; siedzi po prawicy Ojca.}
\w{Et íterum ventúrus est cum glória, judicáre vivos et mórtuos:
   cujus regni non erit finis.}
  {I powtórnie przyjdzie w chwale sądzić żywych i umarłych,
   a królestwu Jego nie będzie końca.}
\w{Et in Spíritum Sanctum, Dóminum, et vivificántem:
   qui ex Patre Filióque procédit.}
  {Wierzę w Ducha Świętego, Pana i Ożywiciela,
   który od Ojca i Syna pochodzi.}
\w{Qui cum Patre et Fílio simul adorátur, et conglorificátur:
   qui locútus est per Prophétas.}
  {Który z Ojcem i Synem wspólnie odbiera uwielbienie i chwałę;
   który mówił przez Proroków.}
\w{Et unam sanctam cathólicam et apostólicam Ecclésiam.}
  {Wierzę w jeden, święty, powszechny i apostolski Kościół.}
\w{Confíteor unum baptísma in remissiónem peccatórum.}
  {Wyznaję jeden chrzest na odpuszczenie grzechów.}
\w{Et exspécto resurrectiónem mortuórum.}
  {I oczekuję wskrzeszenia umarłych.}
\w{Et vitam ventúri saéculi. Amen.}
  {I życia wiecznego w przyszłym świecie. Amen.}
\end{tekst}



% --- IV. Marian antiphons -----------------------------------------------------
\dzial{Antiphonae B. M. V.}{Antyfony maryjne, śpiewane po Mszy świętej według okresu roku liturgicznego}
% Marian antiphons, sung after Mass; \input by every booklet.
% Liturgical year from Advent; each antiphon in the simple, then the solemn tone.
% ---------------------------------------------------------------------------
\antiphonsheading

% ---------------------------------------------------------------------------
\chant{Alma Redemptoris Mater}{tonus V}
{\centering\small\rubr{Od I Niedzieli Adwentu do święta Ofiarowania Pańskiego (2 lutego).}\par}
\par\smallskip{\centering\small\itshape ton prosty\par}
\lilypondfile{../shared/scores/alma-simplex.ly}
\par\smallskip{\centering\small\itshape ton uroczysty\par}
\lilypondfile{../shared/scores/alma.ly}
\begin{tekst}
\w{Alma Redemptóris Mater, quae pérvia caeli porta manes,}
  {Święta Matko Odkupiciela, Bramo niebios zawsze otwarta,}
\w{et stella maris, succúrre cadénti, súrgere qui curat pópulo:}
  {i Gwiazdo morza, wspomóż upadający lud, który pragnie powstać.}
\w{tu quae genuísti, natúra miránte, tuum sanctum Genitórem,}
  {Ty, która ku zdumieniu natury zrodziłaś swego świętego Stwórcę,}
\w{Virgo prius ac postérius, Gabriélis ab ore sumens illud Ave,
   peccatórum miserére.}
  {Dziewico przedtem i potem, z ust Gabriela przyjmująca owo „Ave”,
   zmiłuj się nad grzesznikami.}
\end{tekst}

\clearpage
% ---------------------------------------------------------------------------
\chant{Ave Regina caelorum}{tonus VI}
{\centering\small\rubr{Od Ofiarowania Pańskiego (2 lutego) do Wielkiej Środy.}\par}
\par\smallskip{\centering\small\itshape ton prosty\par}
\lilypondfile{../shared/scores/ave-regina-simplex.ly}
\par\smallskip{\centering\small\itshape ton uroczysty\par}
\lilypondfile{../shared/scores/ave-regina.ly}
\begin{tekst}
\w{Ave, Regína caelórum, ave, Dómina Angelórum:}
  {Witaj, Królowo niebios, witaj, Pani aniołów,}
\w{salve, radix, salve, porta, ex qua mundo lux est orta:}
  {witaj, Korzeniu, witaj, Bramo, z której Światłość wzeszła dla świata.}
\w{Gaude, Virgo gloriósa, super omnes speciósa:}
  {Raduj się, Panno chwalebna, piękniejsza nad wszystkie;}
\w{vale, o valde decóra, et pro nobis Christum exóra.}
  {bądź pozdrowiona, o najpiękniejsza, i przyczyń się za nami do Chrystusa.}
\end{tekst}

\clearpage
% ---------------------------------------------------------------------------
\chant{Regina caeli}{tonus VI}
{\centering\small\rubr{W okresie wielkanocnym, od Niedzieli Zmartwychwstania
do Zesłania Ducha Świętego.}\par}
\par\smallskip{\centering\small\itshape ton prosty\par}
\lilypondfile{../shared/scores/regina-caeli-simplex.ly}
\par\smallskip{\centering\small\itshape ton uroczysty\par}
\lilypondfile{../shared/scores/regina-caeli.ly}
\begin{tekst}
\w{Regína caeli, laetáre, allelúja:}
  {Królowo nieba, wesel się, alleluja,}
\w{quia quem meruísti portáre, allelúja:}
  {bo Ten, któregoś nosić zasłużyła, alleluja,}
\w{resurréxit, sicut dixit, allelúja:}
  {zmartwychwstał, jak powiedział, alleluja.}
\w{ora pro nobis Deum, allelúja.}
  {Módl się za nami do Boga, alleluja.}
\end{tekst}

\clearpage
% ---------------------------------------------------------------------------
\chant{Salve Regina}{tonus V (ton prosty), tonus I (ton uroczysty)}
{\centering\small\rubr{Od uroczystości Najświętszej Trójcy (po Zesłaniu Ducha Świętego)
do Adwentu.}\par}
\par\smallskip{\centering\small\itshape ton prosty\par}
\lilypondfile{../shared/scores/salve-regina-simplex.ly}
\par\smallskip{\centering\small\itshape ton uroczysty\par}
\lilypondfile{../shared/scores/salve-regina.ly}
\begin{tekst}
\w{Salve, Regína, mater misericórdiae;}
  {Witaj, Królowo, Matko miłosierdzia,}
\w{vita, dulcédo et spes nostra, salve.}
  {życie, słodyczy i nadziejo nasza, witaj!}
\w{Ad te clamámus, éxsules fílii Hevae.}
  {Do Ciebie wołamy wygnańcy, synowie Ewy;}
\w{Ad te suspirámus, geméntes et flentes in hac lacrimárum valle.}
  {do Ciebie wzdychamy, jęcząc i płacząc na tym łez padole.}
\w{Eia ergo, advocáta nostra, illos tuos misericórdes óculos
   ad nos convérte.}
  {Przeto, Orędowniczko nasza, one miłosierne oczy Twoje
   na nas zwróć,}
\w{Et Jesum, benedíctum fructum ventris tui,
   nobis post hoc exsílium osténde.}
  {a Jezusa, błogosławiony owoc żywota Twojego,
   po tym wygnaniu nam okaż.}
\w{O clemens, o pia, o dulcis Virgo María.}
  {O łaskawa, o litościwa, o słodka Panno Maryjo!}
\end{tekst}


% --- V. Requiem ----------------------------------------------------------------
\dzial{Missa pro defunctis}{Msza żałobna z obrzędem ostatniego pożegnania}
\renewcommand{\currentmass}{Missa pro defunctis}%  index: chants listed under it
\renewcommand{\masskey}{99}%
\index{Missa 99@Missa pro defunctis}
% Chants of the Missa pro defunctis, the absolution and In paradisum:
% \input by missa-pro-defunctis.tex.  Score paths are relative to a sibling
% directory, so the file also works when \input from another booklet.
% ---------------------------------------------------------------------------
\chant{Requiem aeternam}{introitus · tonus VI}
\lilypondfile{../missa-pro-defunctis/scores/introitus.ly}
\begin{tekst}
\w{Réquiem aetérnam dona eis, Dómine: et lux perpétua lúceat eis.}
  {Wieczny odpoczynek racz im dać, Panie, a światłość wiekuista niechaj im świeci.}
\w{\rubr{Ps. 64.} Te decet hymnus, Deus, in Sion, et tibi reddétur votum in Jerúsalem:
   exáudi oratiónem meam, ad te omnis caro véniet.}
  {\rubr{Ps. 65.} Tobie, Boże, przystoi hymn na Syjonie, Tobie składać się będą śluby
   w Jeruzalem; wysłuchaj mojej modlitwy, do Ciebie przyjdzie wszelkie ciało.}
\w{\rubr{Powtarza się antyfonę} Réquiem aetérnam.}
  {\rubr{Powtarza się antyfonę} Wieczny odpoczynek.}
\end{tekst}

% ---------------------------------------------------------------------------
\chant{Kyrie}{tonus VI}
\lilypondfile{../missa-pro-defunctis/scores/kyrie.ly}
\begin{tekst}
\w{Kýrie, eléison. Christe, eléison. Kýrie, eléison.}
  {Panie, zmiłuj się nad nami. Chryste, zmiłuj się nad nami.
   Panie, zmiłuj się nad nami.}
\end{tekst}

\clearpage
% ---------------------------------------------------------------------------
\chant{Requiem aeternam}{graduale · tonus II}
\lilypondfile{../missa-pro-defunctis/scores/graduale.ly}
\begin{tekst}
\w{Réquiem aetérnam dona eis, Dómine: et lux perpétua lúceat eis.}
  {Wieczny odpoczynek racz im dać, Panie, a światłość wiekuista niechaj im świeci.}
\w{\V In memória aetérna erit justus: ab auditióne mala non timébit.}
  {\V Sprawiedliwy będzie w wiecznej pamięci, nie ulęknie się złej wieści.}
\end{tekst}

% ---------------------------------------------------------------------------
\chant{Absolve, Domine}{tractus · tonus VIII}
\lilypondfile{../missa-pro-defunctis/scores/tractus.ly}
\begin{tekst}
\w{Absólve, Dómine, ánimas ómnium fidélium defunctórum ab omni vínculo delictórum.}
  {Uwolnij, Panie, dusze wszystkich wiernych zmarłych od wszelkich więzów grzechu.}
\w{\V Et grátia tua illis succurrénte, mereántur evádere judícium ultiónis.}
  {\V I niech za wsparciem Twojej łaski zasłużą na uniknięcie wyroku potępienia.}
\w{\V Et lucis aetérnae beatitúdine pérfrui.}
  {\V I niech cieszą się szczęściem wiecznej światłości.}
\end{tekst}

% ---------------------------------------------------------------------------
\chant{Dies irae}{sequentia · tonus I}
\lilypondfile{../missa-pro-defunctis/scores/sequentia.ly}
\begin{tekst}
\w{Dies irae, dies illa, solvet saeclum in favílla: teste David cum Sibýlla.}
  {Dzień gniewu, ów dzień obróci świat w popiół, jak świadczą Dawid i Sybilla.}
\w{Quantus tremor est futúrus, quando judex est ventúrus, cuncta stricte discussúrus!}
  {Jakież będzie drżenie, gdy przyjdzie Sędzia, by wszystko surowo rozsądzić!}
\w{Tuba mirum spargens sonum per sepúlcra regiónum, coget omnes ante thronum.}
  {Trąba, rozsiewając przedziwny dźwięk po grobach wszystkich krain,
   zgromadzi wszystkich przed tronem.}
\w{Mors stupébit et natúra, cum resúrget creatúra, judicánti responsúra.}
  {Zdumieje się śmierć i natura, gdy powstanie stworzenie,
   by dać odpowiedź Sędziemu.}
\w{Liber scriptus proferétur, in quo totum continétur, unde mundus judicétur.}
  {Przyniesiona będzie spisana księga, w której zawarte jest wszystko,
   według czego świat będzie sądzony.}
\w{Judex ergo cum sedébit, quidquid latet apparébit: nil inúltum remanébit.}
  {Gdy więc Sędzia zasiądzie, wyjdzie na jaw wszystko, co ukryte,
   nic nie zostanie bez kary.}
\w{Quid sum miser tunc dictúrus? Quem patrónum rogatúrus, cum vix justus sit secúrus?}
  {Cóż wtedy powiem ja, nędzny? Kogo poproszę o obronę,
   skoro i sprawiedliwy ledwie będzie bezpieczny?}
\w{Rex treméndae majestátis, qui salvándos salvas gratis, salva me, fons pietátis.}
  {Królu straszliwego majestatu, który darmo zbawiasz tych, co mają być zbawieni,
   zbaw mnie, źródło litości.}
\w{Recordáre, Jesu pie, quod sum causa tuae viae: ne me perdas illa die.}
  {Pamiętaj, dobry Jezu, że jestem przyczyną Twojej drogi:
   nie gub mnie w owym dniu.}
\w{Quaerens me, sedísti lassus: redemísti crucem passus: tantus labor non sit cassus.}
  {Szukając mnie, usiadłeś znużony; odkupiłeś mnie, cierpiąc krzyż:
   niech tak wielki trud nie pójdzie na marne.}
\w{Juste judex ultiónis, donum fac remissiónis ante diem ratiónis.}
  {Sprawiedliwy Sędzio, który karzesz, udziel daru przebaczenia
   przed dniem rozrachunku.}
\w{Ingemísco, tamquam reus: culpa rubet vultus meus: supplicánti parce, Deus.}
  {Wzdycham jak winowajca, od winy płonie moje oblicze:
   oszczędź błagającego, Boże.}
\w{Qui Maríam absolvísti, et latrónem exaudísti, mihi quoque spem dedísti.}
  {Ty, który rozgrzeszyłeś Marię i wysłuchałeś łotra,
   i mnie dałeś nadzieję.}
\w{Preces meae non sunt dignae: sed tu bonus fac benígne, ne perénni cremer igne.}
  {Moje prośby nie są godne, lecz Ty, dobry, spraw łaskawie,
   abym nie spłonął w wiecznym ogniu.}
\w{Inter oves locum praesta, et ab haedis me sequéstra, státuens in parte dextra.}
  {Daj mi miejsce wśród owiec i oddziel mnie od kozłów,
   stawiając po prawej stronie.}
\w{Confutátis maledíctis, flammis ácribus addíctis: voca me cum benedíctis.}
  {Gdy potępieni zostaną odrzuceni i wydani srogim płomieniom,
   wezwij mnie z błogosławionymi.}
\w{Oro supplex et acclínis, cor contrítum quasi cinis: gere curam mei finis.}
  {Proszę pokornie, pochylony, z sercem skruszonym jak popiół:
   miej w swej pieczy moją ostatnią godzinę.}
\w{Lacrimósa dies illa, qua resúrget ex favílla judicándus homo reus:
   huic ergo parce, Deus.}
  {Pełen łez będzie ów dzień, w którym powstanie z popiołu
   człowiek winny, by stanąć na sądzie: jemu więc przebacz, Boże.}
\w{Pie Jesu Dómine, dona eis réquiem. Amen.}
  {Dobry Jezu, Panie, daj im odpoczynek. Amen.}
\end{tekst}

% ---------------------------------------------------------------------------
\chant{Domine Jesu Christe}{offertorium · tonus II}
\lilypondfile{../missa-pro-defunctis/scores/offertorium.ly}
\begin{tekst}
\w{Dómine Jesu Christe, Rex glóriae, líbera ánimas ómnium fidélium
   defunctórum de poenis inférni, et de profúndo lacu:}
  {Panie Jezu Chryste, Królu chwały, wybaw dusze wszystkich wiernych
   zmarłych od kar piekła i od otchłani bez dna;}
\w{líbera eas de ore leónis, ne absórbeat eas tártarus, ne cadant in obscúrum:}
  {wybaw je z paszczy lwa, niech ich nie pochłonie piekło,
   niech nie wpadną w ciemności;}
\w{sed sígnifer sanctus Míchael repraeséntet eas in lucem sanctam:}
  {lecz chorąży, święty Michał, niech je stawi w światłości świętej,}
\w{Quam olim Abrahae promisísti, et sémini ejus.}
  {którą niegdyś obiecałeś Abrahamowi i jego potomstwu.}
\w{\V Hóstias et preces tibi, Dómine, laudis offérimus:
   tu súscipe pro animábus illis, quarum hódie memóriam fácimus:
   fac eas, Dómine, de morte transíre ad vitam.}
  {\V Ofiary i modły pochwalne Tobie, Panie, składamy:
   przyjmij je za te dusze, które dziś wspominamy;
   daj im, Panie, przejść od śmierci do życia,}
\w{\rubr{Powtarza się} Quam olim Abrahae \rubr{do końca.}}
  {\rubr{Powtarza się} którą niegdyś obiecałeś \rubr{do końca.}}
\end{tekst}

% ---------------------------------------------------------------------------
\chant{Sanctus}{jak w Mszy XVIII}
\lilypondfile{../missa-pro-defunctis/scores/sanctus.ly}
\begin{tekst}
\w{Sanctus, Sanctus, Sanctus Dóminus Deus Sábaoth.}
  {Święty, Święty, Święty, Pan Bóg Zastępów.}
\w{Pleni sunt caeli et terra glória tua.}
  {Pełne są niebiosa i ziemia chwały Twojej.}
\w{Hosánna in excélsis.}{Hosanna na wysokości.}
\w{Benedíctus qui venit in nómine Dómini.}
  {Błogosławiony, który idzie w imię Pańskie.}
\w{Hosánna in excélsis.}{Hosanna na wysokości.}
\end{tekst}

% ---------------------------------------------------------------------------
\chant{Agnus Dei}{jak w Mszy XVIII}
\lilypondfile{../missa-pro-defunctis/scores/agnus.ly}
\begin{tekst}
\w{Agnus Dei, qui tollis peccáta mundi: dona eis réquiem.}
  {Baranku Boży, który gładzisz grzechy świata, daj im odpoczynek.}
\w{Agnus Dei, qui tollis peccáta mundi: dona eis réquiem.}
  {Baranku Boży, który gładzisz grzechy świata, daj im odpoczynek.}
\w{Agnus Dei, qui tollis peccáta mundi: dona eis réquiem sempitérnam.}
  {Baranku Boży, który gładzisz grzechy świata, daj im wieczny odpoczynek.}
\end{tekst}

\clearpage
% ---------------------------------------------------------------------------
\chant{Lux aeterna}{communio · tonus VIII}
\lilypondfile{../missa-pro-defunctis/scores/communio.ly}
\begin{tekst}
\w{Lux aetérna lúceat eis, Dómine: Cum sanctis tuis in aetérnum, quia pius es.}
  {Światłość wiekuista niechaj im świeci, Panie, ze świętymi Twoimi na wieki,
   bo jesteś łaskawy.}
\w{\V Réquiem aetérnam dona eis, Dómine, et lux perpétua lúceat eis.}
  {\V Wieczny odpoczynek racz im dać, Panie, a światłość wiekuista niechaj im świeci.}
\w{Cum sanctis tuis in aetérnum, quia pius es.}
  {Ze świętymi Twoimi na wieki, bo jesteś łaskawy.}
\end{tekst}

% ---------------------------------------------------------------------------
\chant{Requiescant in pace}{zamiast Ite, missa est}
\lilypondfile{../missa-pro-defunctis/scores/requiescant.ly}
\begin{tekst}
\w{\V Requiéscant in pace.}{\V Niech odpoczywają w pokoju.}
\w{\R Amen.}{\R Amen.}
\end{tekst}

\clearpage
% ---------------------------------------------------------------------------
\begin{center}
  {\LARGE\scshape Absolutio}\par\smallskip
  {\large\itshape Obrzęd ostatniego pożegnania}\par\smallskip
  {\small\rubr{Po Mszy, przy trumnie lub katafalku.}}
\end{center}

\chant{Libera me, Domine}{responsorium · tonus I}
\lilypondfile{../missa-pro-defunctis/scores/libera.ly}
\begin{tekst}
\w{Líbera me, Dómine, de morte aetérna, in die illa treménda:
   Quando caeli movéndi sunt et terra:
   Dum véneris judicáre saeculum per ignem.}
  {Wybaw mnie, Panie, od śmierci wiecznej w ów dzień straszny,
   gdy poruszą się niebiosa i ziemia,
   gdy przyjdziesz sądzić świat przez ogień.}
\w{\V Tremens factus sum ego, et tímeo, dum discússio vénerit, atque ventúra ira.
   Quando caeli movéndi sunt et terra.}
  {\V Drżę cały i lękam się sądu, który nadejdzie, i przyszłego gniewu.
   Gdy poruszą się niebiosa i ziemia.}
\w{\V Dies illa, dies irae, calamitátis et misériae, dies magna et amára valde.
   Dum véneris judicáre saeculum per ignem.}
  {\V Dzień ów, dzień gniewu, klęski i nędzy, dzień wielki i bardzo gorzki.
   Gdy przyjdziesz sądzić świat przez ogień.}
\w{\V Réquiem aetérnam dona eis, Dómine: et lux perpétua lúceat eis.}
  {\V Wieczny odpoczynek racz im dać, Panie, a światłość wiekuista niechaj im świeci.}
\w{\rubr{Powtarza się} Líbera me \rubr{aż do pierwszego wersetu.}}
  {\rubr{Powtarza się} Wybaw mnie \rubr{aż do pierwszego wersetu.}}
\end{tekst}

{\centering\small\rubr{Następnie kapłan mówi: Kýrie, eléison. Christe, eléison. Kýrie,
eléison. Pater noster (po cichu), kropi i okadza trumnę, po czym odmawia modlitwę.}\par}

\clearpage
% ---------------------------------------------------------------------------
\chant{In paradisum}{antiphona · tonus VII}
{\centering\small\rubr{W drodze na cmentarz lub przy wyprowadzeniu ciała z kościoła.}\par}
\lilypondfile{../missa-pro-defunctis/scores/paradisum.ly}
\begin{tekst}
\w{In paradísum dedúcant te Angeli: in tuo advéntu suscípiant te Mártyres,
   et perdúcant te in civitátem sanctam Jerúsalem.}
  {Niech aniołowie zawiodą cię do raju, niech męczennicy wyjdą na twoje spotkanie
   i wprowadzą cię do świętego miasta Jeruzalem.}
\w{Chorus Angelórum te suscípiat, et cum Lázaro quondam páupere aetérnam hábeas réquiem.}
  {Niech cię przyjmie chór aniołów, abyś razem z Łazarzem, niegdyś ubogim,
   miał wieczny odpoczynek.}
\end{tekst}


\vfill
{\centering\itshape + Ad maiorem Dei gloriam +\par}

% --- index --------------------------------------------------------------------
\backmatter
\clearpage
\gdef\headleft{Indeks}\gdef\headright{Indeks}
\phantomsection\addcontentsline{toc}{part}{Indeks}
{\small\printindex}

\end{document}
